%% kapa: exam=1 round=36 year=2025 period=2 form=A subject=appraisal_law subject_ko=감정평가관계법규 seq=39 number=39
%% source: Q-Net 감정평가사 기출문제 (cst00309) artlSeq=5231525 file=raw/1008/5231525/2025년 제36회 감정평가사 1차시험 2교시 문제지 원본.pdf
%% conversion: kordoc 4.24.51 -> tools/gen1.py (auto); review=proofread-sonnet; flags=none
\begin{question}{39}
도시 및 주거환경정비법령상 조합에 관한 설명으로 옳은 것은?
\begin{choices}
  \item 조합원이 정비구역에 위치한 하나의 건축물 또는 토지를 다른 사람과 공유한 경우에도 조합의 임원이 되는데 제한이 없다.
  \item 조합장을 제외한 조합임원은 대의원이 될 수 없다.
  \item 조합장을 제외한 조합임원은 같은 목적의 정비사업을 하는 다른 조합의 임직원을 겸할 수 있다.
  \item 시공자의 선정을 의결하는 총회 및 시공자 선정 취소를 위한 총회의 경우 조합원의 과반수가 직접 출석하여야 한다.
  \item 조합장이 선임 당시 결격사유에 해당하는 자임이 밝혀진 경우 당연 퇴임하고, 퇴임된 임원이 퇴임 전에 관여한 행위는 그 효력을 잃는다.
\end{choices}
\end{question}
\answer{2}
